\ifdefined\gkver\else\def\gkver{student}\fi
\documentclass[\gkver]{gaokaozhenti}
\begin{document}

\chapter{2000年上海}

\begin{enumerate}
\item
%% number: 1
%% paperName: 2000年上海高考第 1 题 5 分
%% typeId: 2
%% type: 多选题
%% chapter: 光学
%% point: 光电效应
%% method: 无
%% score: 5
%% degree: 850
%% duplicateId: 0
%% body:
下列关于光的说法中正确的是\xzanswer{BD}

\fourchoices[answer=BD]
{在真空中红光波长比紫光波长短}
{红光光子能量比紫光光子能量小}
{红光和紫光相遇时能产生干涉现象}
{红光照射某金属时有电子向外发射，紫光照射该金属时一定也有电子向外发射}

%% answer: BD
%% memo:
\memoanswer{本题原卷及材料中未提供详解，待补充。}

\item
%% number: 2
%% paperName: 2000年上海高考第 2 题 5 分
%% typeId: 1
%% type: 单选题
%% chapter: 原子和原子核
%% point: 天然放射性
%% method: 无
%% score: 5
%% degree: 850
%% duplicateId: 0
%% body:
关于 $\alpha$、$\beta$、$\gamma$ 三种射线，下列说法中正确的是\xzanswer{C}

\fourchoices[answer=C]
{$\alpha$ 射线是原子核自发放射出的氦核，它的穿透能力最强}
{$\beta$ 射线是原子核外电子电离形成的电子流，它具有中等的穿透能力}
{$\gamma$ 射线一般伴随着 $\alpha$ 或 $\beta$ 射线产生，它的穿透能力最强}
{$\gamma$ 射线是电磁波，它的穿透能力最弱}

%% answer: C
%% memo:
\memoanswer{本题原卷及材料中未提供详解，待补充。}

\item
%% number: 3
%% paperName: 2000年上海高考第 3 题 5 分
%% typeId: 2
%% type: 多选题
%% chapter: 运动和力
%% point: 牛顿第二定律
%% method: 极值
%% score: 5
%% degree: 650
%% duplicateId: 0
%% body:
一小球用轻绳悬挂在某固定点，现将轻绳水平拉直，然后由静止开始释放小球，考虑小球由静止开始运动到最低位置的过程中\xzanswer{AC}

\fourchoices[answer=AC]
{小球在水平方向的速度逐渐增大}
{小球在竖直方向的速度逐渐增大}
{到达最低位置时小球线速度最大}
{到达最低位置时绳中的拉力等于小球重力}

%% answer: AC
%% memo:
\memoanswer{本题原卷及材料中未提供详解，待补充。}

\item
%% number: 4
%% paperName: 2000年上海高考第 4 题 5 分
%% typeId: 1
%% type: 单选题
%% chapter: 磁场
%% point: 安培力
%% method: 对称法
%% score: 5
%% degree: 650
%% duplicateId: 0
%% body:
如图所示，两根平行放置的长直导线 $a$ 和 $b$ 载有大小相同方向相反的电流，$a$ 受到的磁场力大小为 $F_{1}$，当加入一与导线所在平面垂直的匀强磁场后，$a$ 受到的磁场力大小变为 $F_{2}$，则此时 $b$ 受到的磁场力大小变为\xzanswer{A}

\onepicture[w=3.2cm]{\includesvg{figs/04}}

\fourchoices[answer=A]
{$F_{2}$}
{$F_{1}-F_{2}$}
{$F_{1}+F_{2}$}
{$2F_{1}-F_{2}$}

%% answer: A
%% memo:
\memoanswer{本题原卷及材料中未提供详解，待补充。}

\item
%% number: 5
%% paperName: 2000年上海高考第 5 题 5 分
%% typeId: 2
%% type: 多选题
%% chapter: 功和能
%% point: 动能定理
%% method: 无
%% score: 5
%% degree: 850
%% duplicateId: 0
%% body:
行驶中的汽车制动后滑行一段距离，最后停下；流星在夜空中坠落并发出明亮的光焰；降落伞在空中匀速下降；条形磁铁在下落过程中穿过闭合线圈，线圈中产生电流，上述不同现象中所包含的相同的物理过程是\xzanswer{AD}

\fourchoices[answer=AD]
{物体克服阻力做功}
{物体的动能转化为其它形式的能量}
{物体的势能转化为其它形式的能量}
{物体的机械能转化为其它形式的能量}

%% answer: AD
%% memo:
\memoanswer{本题原卷及材料中未提供详解，待补充。}

\item
%% number: 6
%% paperName: 2000年上海高考第 6 题 5 分
%% typeId: 2
%% type: 多选题
%% chapter: 运动和力
%% point: 变加速直线运动
%% method: 无
%% score: 5
%% degree: 650
%% duplicateId: 0
%% body:
匀速上升的升降机顶部悬挂有一轻质弹簧，弹簧下端挂有一小球，若升降机突然停止，在地面上的观察者看来，小球在继续上升的过程中\xzanswer{AC}

\fourchoices[answer=AC]
{速度逐渐减小}
{速度先增大后减小}
{加速度逐渐增大}
{加速度逐渐减小}

%% answer: AC
%% memo:
\memoanswer{本题原卷及材料中未提供详解，待补充。}

\item
%% number: 7
%% paperName: 2000年上海高考第 7 题 5 分
%% typeId: 2
%% type: 多选题
%% chapter: 机械波
%% point: 波长频率波速
%% method: 无
%% score: 5
%% degree: 650
%% duplicateId: 0
%% body:
如图，沿波的传播方向上有间距均为 $1\Um$ 的六个质点 $a$、$b$、$c$、$d$、$e$、$f$，均静止在各自的平衡位置，一列横波以 $1\Ums$ 的速度水平向右传播，$t=0$ 时到达质点 $a$，$a$ 开始由平衡位置向上运动，$t=1\Us$ 时，质点 $a$ 第一次到达最高点，则在 $4\Us<t<5\Us$ 这段时间内\xzanswer{ACD}

\onepicture[w=7.5cm]{\includesvg{figs/07}}

\fourchoices[answer=ACD]
{质点 $c$ 的加速度逐渐增大}
{质点 $a$ 的速度逐渐增大}
{质点 $d$ 向下运动}
{质点 $f$ 保持静止}

%% answer: ACD
%% memo:
\memoanswer{本题原卷及材料中未提供详解，待补充。}

\item
%% number: 8
%% paperName: 2000年上海高考第 8 题 5 分
%% typeId: 2
%% type: 多选题
%% chapter: 功和能
%% point: 机械能守恒定律
%% method: 整体与隔离
%% score: 5
%% degree: 450
%% duplicateId: 0
%% body:
如图所示，长度相同的三根轻杆构成一个正三角形支架，在 A 处固定质量为 $2m$ 的小球，B 处固定质量为 $m$ 的小球，支架悬挂在 O 点，可绕过 O 点并与支架所在平面相垂直的固定轴转动，开始时 OB 与地面相垂直，放手后开始运动，在不计任何阻力的情况下，下列说法正确的是\xzanswer{BCD}

\onepicture[w=4.0cm]{\includesvg{figs/08}}

\fourchoices[answer=BCD]
{A 球到达最低点时速度为零}
{A 球机械能减少量等于 B 球机械能增加量}
{B 球向左摆动所能达到的最高位置应高于 A 球开始运动时的高度}
{当支架从左向右回摆动时，A 球一定能回到起始高度}

%% answer: BCD
%% memo:
\memoanswer{本题原卷及材料中未提供详解，待补充。}

\item
%% number: 9
%% paperName: 2000年上海高考第 9 题 5 分
%% typeId: 1
%% type: 单选题
%% chapter: 直线运动
%% point: 时间中点速度
%% method: 无
%% score: 5
%% degree: 450
%% duplicateId: 0
%% body:
两木块自左向右运动，现用高速摄影机在同一底片上多次曝光，记录下木块每次曝光时的位置，如图所示，连续两次曝光的时间间隔是相等的。由图可知\xzanswer{C}

\onepicture[w=9.0cm]{\includesvg{figs/09}}

\fourchoices[answer=C]
{在时刻 $t_{2}$ 以及时刻 $t_{5}$ 两木块速度相同}
{在时刻 $t_{1}$ 两木块速度相同}
{在时刻 $t_{3}$ 和时刻 $t_{4}$ 之间某瞬间两木块速度相同}
{在时刻 $t_{4}$ 和时刻 $t_{5}$ 之间某瞬时两木块速度相同}

%% answer: C
%% memo:
\memoanswer{本题原卷及材料中未提供详解，待补充。}

\item
%% number: 10
%% paperName: 2000年上海高考第 10 题 5 分
%% typeId: 2
%% type: 多选题
%% chapter: 电磁感应
%% point: 楞次定律推广
%% method: 无
%% score: 5
%% degree: 450
%% duplicateId: 0
%% body:
如图（a），圆形线圈 P 静止在水平桌面上，其正上方悬挂一相同的线圈 Q，P 和 Q 共轴，Q 中通有变化电流，电流随时间变化的规律如图（b）所示，P 所受的重力为 $G$，桌面对 P 的支持力为 $N$，则\xzanswer{AD}

\onepicture[w=9.0cm, h=3.3cm]{\includesvg{figs/10}}

\fourchoices[answer=AD]
{$t_{1}$ 时刻 $N>G$}
{$t_{2}$ 时刻 $N>G$}
{$t_{3}$ 时刻 $N<G$}
{$t_{4}$ 时刻 $N=G$}

%% answer: AD
%% memo:
\memoanswer{本题原卷及材料中未提供详解，待补充。}

\item
%% number: 11
%% paperName: 2000年上海高考第 11 题 4 分
%% typeId: 3
%% type: 填空题
%% chapter: 电场
%% point: 电势能电势差电场力做功
%% method: 无
%% score: 4
%% degree: 650
%% duplicateId: 0
%% body:
如图，在场强为 $E$ 的匀强电场中有相距为 $l$ 的 A、B 两点，连线 AB 与电场线的夹角为 $\theta$，将一电量为 $q$ 的正电荷从 A 点移到 B 点，若沿直线 AB 移动该电荷，电场力做的功 $W_{1}=$ \tkanswer[2cm]{$qEl\cos\theta$}；若沿路径 ACB 移动该电荷，电场力做的功 $W_{2}=$ \tkanswer[2cm]{$qEl\cos\theta$}；若沿曲线 ADB 移动该电荷，电场力做的功 $W_{3}=$ \tkanswer[2cm]{$qEl\cos\theta$}，由此可知，电荷在电场中移动时，电场力做功的特点是：\tkanswer[5cm]{电场力做功的大小与路径无关，只与始末位置有关。}

\onepicture[w=3.5cm, align=right]{\includesvg{figs/11}}

%% answer:
\jdanswer{$qEl\cos\theta$，$qEl\cos\theta$，$qEl\cos\theta$，电场力做功的大小与路径无关，只与始末位置有关。}

%% memo:
\memoanswer{
$W_{1}=qEl\cos\theta$，$W_{2}=qEl\cos\theta$，$W_{3}=qEl\cos\theta$；由此可知，电场力做功的大小与路径无关，只与始末位置有关。

评分标准：全题 20 分，每小题 4 分，其中第 11 题每格 1 分。
}

\item
%% number: 12
%% paperName: 2000年上海高考第 12 题 4 分
%% typeId: 3
%% type: 填空题
%% chapter: 直线运动
%% point: 匀速直线运动
%% method: 近似与估算
%% score: 4
%% degree: 650
%% duplicateId: 0
%% body:
一架飞机水平匀速地在某同学头顶飞过，当他听到飞机的发动机声从头顶正上方传来时，发现飞机在他前上方约与地面成 $60^{\circ}$ 角的方向上，由此可估算出此飞机的速度约为声速的 \tkanswer[2cm]{$0.58$} 倍。

%% answer:
\jdanswer{$0.58$（或 $\frac{\sqrt{3}}{3}$）}

%% memo:
\memoanswer{
$0.58$（或 $\frac{\sqrt{3}}{3}$）。

评分标准：全题 20 分，每小题 4 分。
}

\item
%% number: 13
%% paperName: 2000年上海高考第 13 题 4 分
%% typeId: 3
%% type: 填空题
%% chapter: 力物体的平衡
%% point: 力矩平衡
%% method: 无
%% score: 4
%% degree: 650
%% duplicateId: 0
%% body:
右图为人手臂面骨骼与肌肉的生理结构示意图，手上托着重量为 $G$ 的物体，

\begin{enumerate}
\item
在方框中画出前臂受力示意图（手、手腕、尺骨和桡骨看成一个整体，所受重力不计，图中 O 点看作固定转动轴，O 点受力可以不画）。
\item
根据图中标尺估算出二头肌此时的收缩力约为 \tkanswer[2cm]{$8G$}。
\end{enumerate}

\onepicture[w=3.5cm, align=right]{figs/13.png}

%% answer:
\jdanswer{
\begin{enumerate}
	\item 
	前臂受力示意图如图所示。
	
	\item
	$8G$
\end{enumerate}
}

%% memo:
\memoanswer{
（1）前臂受力示意图如图所示。

\onepicture[w=3.0cm]{\includesvg{figs/13_an}}

（2）二头肌此时的收缩力约为 $8G$。

评分标准：全题 20 分，每小题 4 分，其中第 13 题受力图 2 分，填空 2 分。
}

\item
%% number: 14
%% paperName: 2000年上海高考第 14 题 4 分
%% typeId: 3
%% type: 填空题
%% chapter: 光学
%% point: 光子说
%% method: 无
%% score: 4
%% degree: 650
%% duplicateId: 0
%% body:
某脉冲激光器的耗电功率为 $2\times10^{3}\UW$，每秒钟输出 10 个光脉冲，每个脉冲持续的时间为 $10^{-8}\Us$，携带的能量为 $0.2\UJ$，则每个脉冲的功率为 \tkanswer[2cm]{$2\times10^{7}$} $\UW$，该激光器将电能转化为激光能量的效率为 \tkanswer[2cm]{$0.001$}。

%% answer:
\jdanswer{$2\times10^{7}$，$0.001$}

%% memo:
\memoanswer{
每个脉冲的功率 $P=\frac{E}{t}=\frac{0.2}{10^{-8}}\UW=2\times10^{7}\UW$；

每秒钟激光器输出的能量 $E_{1}=10\times0.2\UJ=2\UJ$，消耗的电能 $E_{2}=2\times10^{3}\UW\times1\Us=2000\UJ$，则效率 $\eta=\frac{E_{1}}{E_{2}}=\frac{2}{2000}=0.001$。

评分标准：全题 20 分，每小题 4 分，其中第 14 题每格 2 分。
}

\item
%% number: 15
%% paperName: 2000年上海高考第 15 题 4 分
%% typeId: 3
%% type: 填空题
%% chapter: 曲线运动
%% point: 人造地球卫星
%% method: 无
%% score: 4
%% degree: 650
%% duplicateId: 0
%% body:
右图为一名宇航员“漂浮”在地球外层空间的照片，根据照片展现的情景提出两个与物理知识有关的问题（所提的问题可以涉及力学、电磁学、热学、光学、原子物理学等各个部分，只需提出问题，不必作出回答和解释）：例：这名“漂浮”在空中的宇航员相对地球是运动还是静止的？

\begin{enumerate}
\item
\tkanswer[7cm]{此宇航员是否受地球引力作用？}
\item
\tkanswer[7cm]{宇航员背后的天空为什么是黑暗的？}
\end{enumerate}

\onepicture[w=3.5cm, align=right]{figs/15.jpg}

%% answer:
\jdanswer{
\begin{enumerate}
	\item 
	此宇航员是否受地球引力作用？
	
	\item
	宇航员背后的天空为什么是黑暗的？
\end{enumerate}
}

%% memo:
\memoanswer{
只要属于与照片情景有关的物理问题均可得分，例如：此宇航员是否受地球引力作用？此宇航员受力是否平衡？宇航员背后的天空为什么是黑暗的？等等；若解答不以问题的形式出现、不属于物理问题或与照片情景无关，均不能得分。

评分标准：全题 20 分，每小题 4 分，其中第 15 题每格 2 分。
}

\item
%% number: 16
%% paperName: 2000年上海高考第 16 题 4 分
%% typeId: 4
%% type: 实验题
%% chapter: 曲线运动
%% point: 平抛运动
%% method: 无
%% score: 4
%% degree: 650
%% duplicateId: 0
%% body:
右图为用频闪摄影方法拍摄的研究物体作平抛运动规律的照片，图中 A、B、C 为三个同时由同一点出发的小球，AA$'$ 为 A 球在光滑水平面上以速度 $v$ 运动的轨迹；BB$'$ 为 B 球以速度 $v$ 被水平抛出后的运动轨迹；CC$'$ 为 C 球自由下落的运动轨迹，通过分析上述三条轨迹可得出结论：\tkanswer[6cm]{作平抛运动的物体在水平方向作匀速直线运动，在竖直方向作自由落体运动}。

\onepicture[w=3.0cm, align=right]{figs/16.jpg}

%% answer:
\jdanswer{平抛运动是水平方向的匀速直线运动和竖直方向的自由落体运动的合成。}

%% memo:
\memoanswer{
作平抛运动的物体在水平方向作匀速直线运动，在竖直方向作自由落体运动（或平抛运动是水平方向的匀速直线运动和竖直方向的自由落体运动的合成）。（4 分）

评分标准：全题 30 分，具体见于各题答案后。
}

\item
%% number: 17
%% paperName: 2000年上海高考第 17 题 4 分
%% typeId: 4
%% type: 实验题
%% chapter: 光学
%% point: 光的衍射
%% method: 无
%% score: 4
%% degree: 650
%% duplicateId: 0
%% body:
单色光源发出的光经一狭缝，照射到光屏上，可观察到的图象是\xzanswer{A}

\fourchoices[answer=A, ispicture=true, h=0.9cm]
{figs/17a.png}
{figs/17b.png}
{figs/17c.png}
{figs/17d.png}

%% answer:
\jdanswer{A}

%% memo:
\memoanswer{
单色光经单缝后发生衍射，在光屏上形成中央亮纹最宽最亮、两侧亮纹逐渐变窄变暗的图样，故选 A。（4 分）

评分标准：全题 30 分，具体见于各题答案后。
}

\item
%% number: 18
%% paperName: 2000年上海高考第 18 题 7 分
%% typeId: 4
%% type: 实验题
%% chapter: 力物体的平衡
%% point: 力矩盘
%% method: 无
%% score: 7
%% degree: 650
%% duplicateId: 0
%% body:
用右图所示装置做“研究有固定转动轴物体平衡条件”的实验，力矩盘上各同心圆的间距相等。

\begin{enumerate}
\item
在用细线悬挂钩码前，以下哪些措施是必要的\xzanswer{ACD}

\fourchoices[answer=ACD]
{判断力矩盘是否处在竖直平面}
{判断横杆 $MN$ 是否严格保持水平}
{判断力矩盘与转轴间的摩擦是否足够小}
{判断力矩盘的重心是否位于盘中心}
\item
在 A、B、C 三点分别用细线悬挂钩码后，力矩盘平衡，如图所示，已知每个钩码所受的重力为 $1\UN$，则此时弹簧秤示数为 \tkanswer[2cm]{$4\UN$}。
\item
现同时撤去 B、C 两处的钩码，而改用一根细线悬挂 $5\UN$ 的钩码，为使力矩盘仍然保持原平衡状态，且弹簧秤示数不变，试在图中用直线画出该细线在力矩盘上的悬挂位置。
\end{enumerate}

\onepicture[h=4.5cm, align=right]{\includesvg{figs/18}}

%% answer:
\jdanswer{
\begin{enumerate}
	\item 
	ACD
	
	\item
	$4\UN$
	
	\item
	悬线在转轴右侧竖直方向与第二圆环相切，如图所示。
\end{enumerate}
}

%% memo:
\memoanswer{
（1）在用细线悬挂钩码前，应判断力矩盘是否处在竖直平面、与转轴间的摩擦是否足够小、重心是否位于盘中心，故选 ACD。（4 分）

（2）由力矩平衡条件可得此时弹簧秤示数为 $4\UN$。（2 分）

（3）撤去 B、C 两处钩码，改用一根细线悬挂 $5\UN$ 的钩码，且弹簧秤示数不变，则悬线在转轴右侧竖直方向与第二圆环相切，如图所示（钩码不画不扣分）。（1 分）

\onepicture[h=4.0cm]{\includesvg{figs/18_an}}

评分标准：全题 30 分，具体见于各题答案后，其中第 18（1）有多个正确选项，全选对得 4 分，选对但不全得 2 分，有选错或全部不选得 0 分。
}

\item
%% number: 19
%% paperName: 2000年上海高考第 19 题 8 分
%% typeId: 4
%% type: 实验题
%% chapter: 电路
%% point: 故障分析
%% method: 无
%% score: 8
%% degree: 650
%% duplicateId: 0
%% body:
某同学按如图所示电路进行实验，实验时该同学将变阻器的触片 P 移到不同位置时测得各电表的示数如下表所示。

\begin{table}[htbp]
\centering
\begin{tblr}{|c|c|c|c|c|}
\hline
序号 & $A_{1}$ 示数（$\UA[a]$） & $A_{2}$ 示数（$\UA[a]$） & $V_{1}$ 示数（$\UV[a]$） & $V_{2}$ 示数（$\UV[a]$） \\
\hline
1 & 0.60 & 0.30 & 2.40 & 1.20 \\
\hline
2 & 0.44 & 0.32 & 2.56 & 0.48 \\
\hline
\end{tblr}
\end{table}

将电压表内阻看作无限大，电流表内阻看作零。

\begin{enumerate}
\item
电路中 $E$、$r$ 分别为电源的电动势和内阻，$R_{1}$、$R_{2}$、$R_{3}$ 为定值电阻，在这五个物理量中，可根据上表中的数据求得的物理量是（不要求具体计算）\tkanswer[3cm]{$R_{2}$、$R_{3}$、$E$}。
\item
由于电路发生故障，发现两电压表示数相同了（但不为零），若这种情况的发生是由用电器引起的，则可能的故障原因是 \tkanswer[4cm]{$R_{p}$ 短路或 $R_{2}$ 断开}。
\end{enumerate}

\onepicture[h=4.5cm, align=right]{\includesvg{figs/19}}

%% answer:
\jdanswer{
\begin{enumerate}
	\item 
	$R_{2}$、$R_{3}$、$E$
	
	\item
	$R_{p}$ 短路或 $R_{2}$ 断开
\end{enumerate}
}

%% memo:
\memoanswer{
（1）由表中数据可求得 $R_{3}=\frac{U_{1}}{I_{2}}$、$R_{2}=\frac{U_{2}}{I_{1}-I_{2}}$，并由 $E=U_{1}+I_{1}(R_{1}+r)$ 联立两组数据可求得电源电动势 $E$，故可求得 $R_{2}$、$R_{3}$、$E$。（1 分，1 分，2 分）

（2）两电压表示数相同且不为零，可能的故障原因是变阻器 $R_{p}$ 短路或 $R_{2}$ 断开。（4 分）

评分标准：全题 30 分，具体见于各题答案后。19（1）若填 $r$、$R_{1}$ 的，在本小题的得分中，多填一个扣 1 分，扣完为止，如填 $R_{p}$，不扣分。
}

\item
%% number: 20
%% paperName: 2000年上海高考第 20 题 7 分
%% typeId: 4
%% type: 实验题
%% chapter: 气体性质
%% point: 盖吕萨克定律
%% method: 无
%% score: 7
%% degree: 450
%% duplicateId: 0
%% body:
现有一根粗细均匀长约 $40\Ucm$，两端开口的玻璃管，内有一小段水柱，一个弹簧秤，一把毫米刻度尺，一小块橡皮泥，一个足够高的玻璃容器，内盛有冰和水的混合物，选用合适的器材，设计一个实验，估测当时的室内温度，要求：

\begin{enumerate}
\item
在右边方框中画出实验示意图；
\item
写出要测定的物理量 \tkanswer[3cm]{室温中空气柱的长度 $l_{1}$、冰水中空气柱的长度 $l_{2}$}，写出可直接应用的物理量 \tkanswer[2cm]{冰水的温度 $T_{0}$}；
\item
写出计算室温的表达式 \tkanswer[2cm]{$T=\frac{l_{1}}{l_{2}}T_{0}$}。
\end{enumerate}

\onepicture[w=4.5cm, align=right]{figs/20.png}

%% answer:
\jdanswer{
\begin{enumerate}
	\item 
	实验示意图如图所示。
	
	\item
	$l_{1}$、$l_{2}$、$T_{0}$
	
	\item
	$T=\frac{l_{1}}{l_{2}}T_{0}$
\end{enumerate}
}

%% memo:
\memoanswer{
（1）实验示意图如图所示（玻璃管用水柱密封，分别放在室温空气和冰水混合物中）。（3 分）

\onepicture[w=3.5cm]{figs/20_an.png}

（2）要测定的物理量：玻璃管放在室温中时空气柱的长度 $l_{1}$、玻璃管浸在冰水内时空气柱的长度 $l_{2}$；可直接应用的物理量：冰水的温度 $T_{0}$。（2 分，每个空格 1 分）

（3）由等压变化 $\frac{l_{1}}{T}=\frac{l_{2}}{T_{0}}$ 可得计算室温的表达式 $T=\frac{l_{1}}{l_{2}}T_{0}$。（2 分）

评分标准：全题 30 分，具体见于各题答案后。20（1）中，若画入弹簧秤，其余正确，扣 1 分，玻璃管不密封的不能得分。20（2）中 $T_{0}$ 用具体数值表示的，同样给分。
}

\item
%% number: 21
%% paperName: 2000年上海高考第 21 题 12 分
%% typeId: 6
%% type: 计算题
%% chapter: 运动和力
%% point: 牛顿定律的应用
%% method: 无
%% score: 12
%% degree: 650
%% duplicateId: 0
%% body:
风洞实验室中可以产生水平方向的、大小可调节的风力，现将一套有小球的细直杆放入风洞实验室，小球孔径略大于细杆直径。

\begin{enumerate}
\item
当杆在水平方向上固定时，调节风力的大小，使小球在杆上作匀速运动，这时小球所受的风力为小球所受重力的 0.5 倍，求小球与杆间的滑动摩擦因数。
\item
保持小球所受风力不变，使杆与水平方向间夹角为 $37^{\circ}$ 并固定，则小球从静止出发在细杆上滑下距离 $s$ 所需时间为多少？（$\sin37^{\circ}=0.6$，$\cos37^{\circ}=0.8$）
\end{enumerate}

\onepicture[w=5.5cm, align=right]{\includesvg{figs/21}}

%% answer:
\jdanswer{
\begin{enumerate}
	\item 
	$0.5$
	
	\item
	$\sqrt{\frac{8s}{3g}}$
\end{enumerate}
}

%% memo:
\memoanswer{
（1）设小球所受的风力为 $F$，小球质量为 $m$。

$F=\mu mg$ ①

$\mu=\frac{F}{mg}=\frac{0.5mg}{mg}=0.5$ ②

（2）设杆对小球的支持力为 $N$，摩擦力为 $f$。

沿杆方向 $F\cos\theta+mg\sin\theta-f=ma$ ③

垂直于杆方向 $N+F\sin\theta-mg\cos\theta=0$ ④

$f=\mu N$ ⑤

可解得

$a=\frac{F\cos\theta+mg\sin\theta-f}{m}=\left(g+\frac{F^{2}}{m^{2}g}\right)\sin\theta=\frac{3}{4}g$ ⑥

$s=\frac{1}{2}at^{2}$ ⑦

$t=\sqrt{\frac{2s}{\frac{3}{4}g}}=\sqrt{\frac{8s}{3g}}$ ⑧

\onepicture[w=4.5cm]{\includesvg{figs/21_an}}

评分标准：全题 12 分。（1）3 分。正确得出②式，得 3 分。仅写出①式，得 1 分。（2）9 分，正确得出⑥式，得 6 分，仅写出③、④式，各得 2 分，仅写出⑤式，得 1 分，正确得出⑧式，得 3 分，仅写出⑦式，得 2 分，$g$ 用数值代入的不扣分。
}

\item
%% number: 22
%% paperName: 2000年上海高考第 22 题 12 分
%% typeId: 6
%% type: 计算题
%% chapter: 气体性质
%% point: 气体多过程
%% method: 临界
%% score: 12
%% degree: 450
%% duplicateId: 0
%% body:
如图所示，粗细均匀，两端开口的 U 形管竖直放置，管的内径很小，水平部分 BC 长 $14\Ucm$，一空气柱将管内水银分隔成左右两段，大气压强相当于高为 $76\UcmHg$ 的压强。

\begin{enumerate}
\item
当空气柱温度为 $T_{0}=273\UK$，长为 $l_{0}=8\Ucm$ 时，BC 管内左边水银柱长 $2\Ucm$，AB 管内水银柱长也是 $2\Ucm$，则右边水银柱总长是多少？
\item
当空气柱温度升高到多少时，左边的水银恰好全部进入竖直管 AB 内？
\item
当空气柱温度为 $490\UK$ 时，两竖直管内水银柱上表面高度各为多少？
\end{enumerate}

\onepicture[w=4.0cm, align=right]{\includesvg{figs/22}}

%% answer:
\jdanswer{
\begin{enumerate}
	\item 
	$6\Ucm$
	
	\item
	$420\UK$
	
	\item
	右管内水银面高度为 $4\Ucm$，左管内水银上表面高度为 $6\Ucm$
\end{enumerate}
}

%% memo:
\memoanswer{
（1）U 形管两端均开口，所以两竖直管内水银面高度应相同，即右边竖直管内水银柱高度为

$h_{0}=2\Ucm$ ①

右边水平管内水银柱长度为

$14-l_{0}-2=4\Ucm$

右边水银柱总长是

$4+2=6\Ucm$ ②

（2）左边的水银全部进入竖直管内时，两竖直管内水银面高度均为

$h_{1}=4\Ucm$ ③

此时，右边水平管内水银柱长度为 $2\Ucm$，所以空气柱长为

$l_{1}=14-2=12\Ucm$ ④

$\frac{(p_{0}+h_{0})l_{0}}{T_{0}}=\frac{(p_{0}+h_{1})l_{1}}{T_{1}}$ ⑤

$T_{1}=\frac{(p_{0}+h_{1})l_{1}}{(p_{0}+h_{0})l_{0}}T_{0}=\frac{80\times12}{78\times8}\times273\UK=420\UK$ ⑥

（3）设温度为 $T_{2}=490\UK$ 时，空气柱长为 $l_{2}$。

等压过程 $\frac{l_{1}}{T_{1}}=\frac{l_{2}}{T_{2}}$ ⑦

$l_{2}=\frac{l_{1}}{T_{1}}T_{2}=\frac{12}{420}\times490\Ucm=14\Ucm$ ⑧

其中有 $2\Ucm$ 进入左边竖直管内，故

右管内水银面高度为 $h_{1}=4\Ucm$ ⑨

左管内水银上表面高度为 $h_{2}=4+2=6\Ucm$ ⑩

评分标准：全题 12 分。（1）3 分，正确得出结果②，得 3 分，仅得出①式，得 1 分。（2）5 分，正确得出④式，得 2 分，仅得出③式，得 1 分，正确得出⑥式，得 3 分，仅得出⑤式，得 2 分。（3）4 分，正确得出⑧式，得 2 分，仅写出⑦式，得 1 分，得出⑨式、⑩式，各 1 分。
}

\item
%% number: 23
%% paperName: 2000年上海高考第 23 题 13 分
%% typeId: 6
%% type: 计算题
%% chapter: 电磁感应
%% point: 电磁感应受力分析
%% method: 无
%% score: 13
%% degree: 450
%% duplicateId: 0
%% body:
如图所示，固定水平桌面上的金属框架 cdef，处在竖直向下的匀强磁场中，金属棒 ab 搁在框架上，可无摩擦滑动，此时 adeb 构成一个边长为 $l$ 的正方形，棒的电阻为 $r$，其余部分电阻不计，开始时磁感强度为 $B_{0}$。

\begin{enumerate}
\item
若从 $t=0$ 时刻起，磁感强度均匀增加，每秒增量为 $k$，同时保持棒静止，求棒中的感应电流，在图上标出感应电流的方向。
\item
在上述（1）情况中，始终保持棒静止，当 $t=t_{1}$ 秒末时需加的垂直于棒的水平拉力为多大？
\item
若从 $t=0$ 时刻起，磁感强度逐渐减小，当棒以恒定速度 $v$ 向右作匀速运动时，可使棒中不产生感应电流，则磁感强度应怎样随时间变化（写出 $B$ 与 $t$ 的关系式）？
\end{enumerate}

\onepicture[w=6.0cm, align=right]{\includesvg{figs/23}}

%% answer:
\jdanswer{
\begin{enumerate}
	\item 
	$I=\frac{kl^{2}}{r}$，方向：逆时针（如图所示）
	
	\item
	$F=(B_{0}+kt_{1})\frac{kl^{3}}{r}$
	
	\item
	$B=\frac{B_{0}l}{l+vt}$
\end{enumerate}
}

%% memo:
\memoanswer{
（1）感应电动势 $\varepsilon=\frac{\Delta\Phi}{\Delta t}=kl^{2}$ ①

感应电流 $I=\frac{\varepsilon}{r}=\frac{kl^{2}}{r}$ ②

方向：逆时针（如图）。

\onepicture[w=5.0cm]{\includesvg{figs/23_an}}

（2）$t=t_{1}$ 秒时，$B=B_{0}+kt_{1}$ ③

$F=BIL$ ④

$F=(B_{0}+kt_{1})\frac{kl^{3}}{r}$ ⑤

（3）总磁通量不变 $Bl(l+vt)=B_{0}l^{2}$ ⑥

$B=\frac{B_{0}l}{l+vt}$ ⑦

评分标准：全题 13 分。（1）5 分，正确得出②式，得 4 分，仅得出①式得 2 分，方向正确，得 1 分。（2）4 分，正确得出⑤式，得 4 分，仅得出③式，得 1 分，仅写出④式，得 2 分。（3）4 分，正确得出⑦式，得 4 分，仅得出⑥式，得 2 分。
}

\item
%% number: 24
%% paperName: 2000年上海高考第 24 题 13 分
%% typeId: 6
%% type: 计算题
%% chapter: 光学
%% point: 光的电磁说
%% method: 建立模型
%% score: 13
%% degree: 250
%% duplicateId: 0
%% body:
阅读如下资料并回答问题：

自然界中的物体由于具有一定的温度，会不断向外辐射电磁波，这种辐射因与温度有关，称为热辐射，热辐射具有如下特点：

①辐射的能量中包含各种波长的电磁波；

②物体温度越高，单位时间从物体表面单位面积上辐射的能量越大；

③在辐射的总能量中，各种波长所占的百分比不同。

处于一定温度的物体在向外辐射电磁能量的同时，也要吸收由其他物体辐射的电磁能量，如果它处在平衡状态，则能量保持不变，若不考虑物体表面性质对辐射与吸收的影响，我们定义一种理想的物体，它能 $100\%$ 地吸收入射到其表面的电磁辐射，这样的物体称为黑体，单位时间内从黑体表面单位面积辐射的电磁波的总能量与黑体绝对温度的四次方成正比，即 $P_{0}=\sigma T^{4}$，其中常量 $\sigma=5.67\times10^{-8}\UW/(\Umq\cdot\UK^{4})$。

在下面的问题中，把研究对象都简单地看作黑体。

有关数据及数学公式：太阳半径 $R_{s}=696\,000\Ukm$，太阳表面温度 $T=5770\UK$，火星半径 $r=3395\Ukm$，球面积 $S=4\pi R^{2}$，其中 $R$ 为球半径。

\begin{enumerate}
\item
太阳热辐射能量的绝大多数集中在波长为 $2\times10^{-7}\Um\sim1\times10^{-5}\Um$ 范围内，求相应的频率范围。
\item
每小时从太阳表面辐射的总能量为多少？
\item
火星受到来自太阳的辐射可认为垂直射到面积为 $\pi r^{2}$（$r$ 为火星半径）的圆盘上，已知太阳到火星的距离约为太阳半径的 400 倍，忽略其它天体及宇宙空间的辐射，试估算火星的平均温度。
\end{enumerate}

%% answer:
\jdanswer{
\begin{enumerate}
	\item 
	$3\times10^{13}\UHz\sim1.5\times10^{15}\UHz$
	
	\item
	$W=1.38\times10^{30}\UJ$
	
	\item
	$T'=204\UK$
\end{enumerate}
}

%% memo:
\memoanswer{
（1）$f=\frac{c}{\lambda}$ ①

$f_{1}=\frac{3.00\times10^{8}}{2\times10^{-7}}\UHz=1.5\times10^{15}\UHz$ ②

$f_{2}=\frac{3.00\times10^{8}}{1\times10^{-5}}\UHz=3\times10^{13}\UHz$ ③

辐射的频率范围为 $3\times10^{13}\UHz\sim1.5\times10^{15}\UHz$。

（2）每小时从太阳表面辐射的总能量为

$W=4\pi\sigma R_{s}^{2}T^{4}t$ ④

代入数据得 $W=1.38\times10^{30}\UJ$。 ⑤

（3）设火星表面温度为 $T'$，太阳到火星距离为 $d$，火星单位时间内吸收来自太阳的辐射能量为

$P_{\text{in}}=4\pi\sigma R_{s}^{2}T^{4}\cdot\frac{\pi r^{2}}{4\pi d^{2}}$ ⑥

$d=400R_{s}$

$P_{\text{in}}=\frac{\pi\sigma T^{4}r^{2}}{400^{2}}$ ⑦

火星单位时间内向外辐射电磁波能量为

$P_{\text{out}}=4\pi\sigma r^{2}T'^{4}$ ⑧

火星处在平衡状态 $P_{\text{in}}=P_{\text{out}}$ ⑨

即 $\frac{\pi\sigma T^{4}r^{2}}{400^{2}}=4\pi\sigma r^{2}T'^{4}$ ⑩

由⑩式解得火星平均温度 $T'=\frac{T}{\sqrt{800}}=204\UK$。（11）

评分标准：全题 13 分。（1）3 分，正确得出①、②、③式，各得 1 分。（2）5 分，正确得出⑤式，得 5 分，仅写出④式，得 3 分。（3）5 分，正确得出⑩式，得 4 分，仅写出⑥式或⑦式，得 1 分；仅写出⑧式，得 1 分，正确得出（11）式，得 1 分。
}

\end{enumerate}

\end{document}
